% problem-000199 2 氯化铵铜复盐组成
\section*{2. 氯化铵铜复盐组成}
在 $\mathrm { N H _ { 4 } C l  – C u C l _ { 2 } – H _ { 2 } O }$ 体系中，结晶出⼀种淡蓝⾊的物质A，其组成可表⽰为 $x \mathrm { N H _ { 4 } C l { \cdot } C u C l _ { 2 } { \cdot } } y \mathrm { H _ { 2 } O _ { \circ } }$ 称取1.4026 g 晶体A，溶于⽔并在250 mL容量瓶中定容。Cl−分析：移取25.00 mL溶液，加⼊2 滴0.5%荧光⻩的⼄醇溶液和⼀滴稀 NaOH 溶液，再加 2 mL 0.5%淀粉溶液，⽤ 0.1036 mol L−1 $\mathrm { A g N O _ { 3 } }$ 溶液滴定（反应1）⾄出现粉红⾊，消耗19.52 mL； $\mathrm { C u ^ { 2 + } }$ 分析：移取 25.00 mL 溶液，加⼊ 1 mol L−1 H $\mathrm { _ { 2 S O _ { 4 } } }$ 溶液5mL，再加固体 KI 1.5 g，混匀并放置（反应 2），⽤ 0.02864 mol L−1 $\mathrm { N a _ { 2 } S _ { 2 } O _ { 3 } }$ 溶液滴定（反应3）⾄棕⾊较浅时加⼊2mL0.5%淀粉溶液，继续滴⾄蓝紫⾊恰好消失，消耗17.65 mL。

\subsection*{2-1}
写出反应1-3的⽅程式。

\subsection*{2-2}
通过计算确定 A 中�和�的值（NH_{4}Cl 式量 53.49， $\mathrm { C u C l _ { 2 } }$ 式量 134.5）。
